\chapter{Liver-Kidney Microsomal Antibodies}

\section*{What Is This Test?}
Liver-kidney microsomal type 1 (LKM-1) antibody testing detects autoantibodies, usually directed against cytochrome P450 2D6, that are associated with autoimmune hepatitis type 2. This form is more frequent in children and young adults, but it can occur at any age.

\section*{Alternative Names}
Anti-LKM-1; Liver-Kidney Microsomal Antibody; LKM Antibody; Autoimmune Hepatitis Type 2 Antibody

\section*{How It Is Performed}
LKM antibodies are measured in serum by indirect immunofluorescence or antigen-specific immunoassay. Testing is usually combined with antinuclear antibody, smooth-muscle antibody, immunoglobulin G, and other liver autoantibodies as clinically indicated.

\section*{Why It Is Ordered}
\begin{itemize}
  \item Evaluation of unexplained hepatocellular liver-test abnormalities when autoimmune hepatitis is suspected
  \item Investigation of suspected autoimmune hepatitis in children or young adults
  \item Classification of autoimmune hepatitis when routine ANA and smooth-muscle antibody tests are negative
\end{itemize}

\section*{Interpretation and Limitations}
A positive result supports autoimmune hepatitis type 2 but is not diagnostic by itself. Low-level reactivity can occur in other liver diseases and infections, while a negative result does not exclude autoimmune hepatitis. Interpretation requires liver enzymes, IgG, other autoantibodies, exclusion of viral and drug-related causes, and specialist assessment; liver biopsy is sometimes used for confirmation but is a procedure outside this test chapter.
\section*{Specimen and Preparation}
The test uses serum and usually requires no fasting. Tell the clinician about liver medicines, supplements, viral hepatitis risk, alcohol use, and other autoimmune disease. The assay method and titre matter because indirect immunofluorescence and antigen-specific methods may use different reporting conventions.

\section*{Risks and Safety}
Physical risk is limited to venipuncture. A positive antibody is not, by itself, proof of autoimmune hepatitis and should not lead to immunosuppressive treatment without complete clinical assessment. Jaundice, confusion, bleeding, severe abdominal swelling, or rapidly worsening illness requires urgent review.

\section*{Understanding Results and Follow-up}
A positive LKM-1 result supports autoimmune hepatitis type 2 in the right context; a negative result does not exclude autoimmune hepatitis. Follow-up commonly includes ALT, AST, bilirubin, alkaline phosphatase, IgG, other autoantibodies, viral testing, and specialist assessment. Interpretation must account for drug-induced and infectious liver disease.

\section*{Regulatory Status and Authorization Date}
This chapter describes a test category, not a specific commercial product. FDA, CDSCO, EU IVDR, and MHRA approval or registration dates therefore cannot be assigned without the assay manufacturer, product name, instrument, and intended use. For laboratory-developed tests, an individual agency approval date may not exist. See the Preface for the interpretation of regulatory status.

\clearpage
